\chapter{Introduction}
\label{ch:intro}

\section{Background}

A connected autonomous vehicle (CAV) is in continuous radio contact with the
infrastructure around it. It broadcasts its position, speed and heading about
ten times a second, reports hazards such as a collision ahead or ice on the
road, and receives warnings issued by roadside units (RSUs) and edge servers.
Every one of these messages is safety-relevant: a forged ``brake now'' or a
suppressed hazard warning can cause a crash. Vehicle-to-everything (V2X)
security therefore rests on two guarantees, that a message really comes from
the registered entity that claims to have sent it, and that it has not been
altered or replayed.

Deployed V2X security, specified in IEEE~1609.2~\cite{ieee16092}, provides
these guarantees with elliptic-curve cryptography: ECDSA signatures on
messages and certificates, and ECDH where confidentiality is needed. Both
rest on the elliptic-curve discrete-logarithm problem. Shor's
algorithm~\cite{shor} solves that problem, and integer factorisation, in
polynomial time on a sufficiently large quantum computer. Such machines do not
exist yet, but two facts make the problem urgent now. First, an adversary can
record encrypted traffic today and decrypt it once the hardware arrives
(``harvest now, decrypt later''~\cite{mosca}); a vehicle's movement history
stays sensitive for years. Second, vehicles stay in service for fifteen years
or more, so the cryptography fitted to a car sold today must still be secure
in 2040.

In August 2024 NIST published the first post-quantum standards: ML-KEM
(FIPS~203~\cite{fips203}), a key-encapsulation mechanism derived from
CRYSTALS-Kyber, and ML-DSA (FIPS~204~\cite{fips204}), a signature scheme
derived from CRYSTALS-Dilithium. Both are built on the Module Learning With
Errors problem over lattices, for which no efficient quantum algorithm is
known, and both run on ordinary processors. Their cost is size: an ML-DSA-65
public key is 1,952 bytes and a signature 3,309 bytes, against 65 and about
72 bytes for ECDSA P-256. On a 6\,Mb/s IEEE 802.11p channel shared by every
vehicle in range, that difference matters more than computation time.

\section{Problem statement}

Our base paper, Aslam, Bhardwaj and Chaudhary~\cite{base}, proposes a
quantum-resilient V2X framework: a Trusted Authority registers entities,
vehicles establish keys with RSUs using Kyber, edge servers route events by
criticality, and an Adaptive Grouping Score PBFT (AGS-PBFT) consensus layer
records critical events on a ledger. The paper reports latency and message
counts but leaves open the questions a deployment depends on:

\begin{enumerate}
\item Can the framework be implemented as specified with the standardised
algorithms, and is the resulting protocol secure?
\item Do post-quantum handshakes and safety messages fit vehicular latency
budgets (100\,ms for periodic beacons) when computation \emph{and} radio
transmission are both measured?
\item How does a shared 802.11p channel behave as vehicle density grows, when
every handshake carries post-quantum keys, signatures and certificates?
\item Does AGS-PBFT deliver the savings the paper claims over PBFT, and at
what cost?
\end{enumerate}

\section{Objectives}

\begin{enumerate}
\item Implement the complete framework (Trusted Authority, registration,
mutually authenticated ML-KEM/ML-DSA handshake, record layer, edge routing,
AGS-PBFT with view change, and a ledger) as a working distributed system.
\item Correct every defect in the published design that prevents it from
working or makes it insecure, and document each correction.
\item Measure computation, message size and radio airtime for all NIST
security levels and for a classical ECDH/ECDSA baseline run through the same
protocol code.
\item Evaluate the protocol in ns-3 with IEEE 802.11p, fading and vehicle
mobility from 50 to 200 vehicles.
\item Compare AGS-PBFT with PBFT using real signatures, and test both under
Byzantine faults and leader crashes.
\item Test the implementation against concrete attacks.
\item Reduce the handshake's size, the dominant cost, with certificate caching.
\end{enumerate}

\section{Contributions}

\begin{itemize}
\item \textbf{A complete, running implementation} (the \code{qrcav} package,
about 7,600 lines of Python, including tests, and 870 lines of C++ for ns-3) that runs as 19
processes over TCP, inside ns-3, and as an evaluation suite, with 57 automated
tests.
\item \textbf{Corrections to the base paper} (D1--D5): a consensus layer that
signs with a KEM, a circular key-derivation listing, a session key that
hashes a value with itself, key escrow at the Trusted Authority, and a ledger
platform that cannot host PBFT. \textbf{Corrections to our own Review~1
design} (F1--F7), the most serious being the absence of client
authentication.
\item \textbf{Design decisions where the paper is undefined} (D6--D9) that
make AGS-PBFT implementable: regrouping when all scores are equal, how
candidates earn promotion, how group size follows load, and how the leader's
own vote counts.
\item \textbf{Certificate caching} for the post-quantum handshake, with its
security argument, which halves the handshake to \V{deck-cached-bytes} bytes.
\item \textbf{A measured answer} to the questions above: per-message
post-quantum protection fits the budget; channel capacity during handshakes
does not scale; the cause is size, which certificate caching and
exponential backoff after failed handshakes reduce.
\end{itemize}

\section{Scope}

The project uses ordinary (classical) hardware; post-quantum cryptography is
software and needs no quantum devices. Out of scope are deployment on real
on-board units, side-channel hardening beyond what the OpenSSL implementation
provides, formal verification of the protocol, and pseudonymous certificates
for location privacy.

\section{Organisation of the report}

Chapter~\ref{ch:lit} surveys related work. Chapter~\ref{ch:bg} explains the
cryptographic and consensus building blocks. Chapter~\ref{ch:design}
presents the system design, the defects found and the corrections.
Chapter~\ref{ch:impl} describes the implementation and the simulation
methodology. Chapter~\ref{ch:results} reports and discusses the results, and
Chapter~\ref{ch:conc} concludes with limitations and future work.
